\documentclass[10pt]{article}
\usepackage[margin=20mm]{geometry}
\usepackage{amsmath,amssymb,fontspec,unicode-math,xurl,hyperref}
\IfFontExistsTF{DejaVu Serif}{\setmainfont{DejaVu Serif}}{\setmainfont{Latin Modern Roman}}
\setmathfont{latinmodern-math.otf}
\setlength{\parindent}{0pt}\setlength{\parskip}{6pt}\sloppy
\title{Technical Verification: On Linear Convergence in Smooth Convex-Concave Bilinearly-Coupled Saddle-Point Optimization: Lower Bounds and Optimal Algorithms}
\author{AI-assisted technical review}
\date{}
\begin{document}
\maketitle
Abstract: A quantifier problem is established in the proof of the all-time distance lower bound. The proposed finite-dimensional witness admits an allowed gradient method whose distance to its solution tends to zero. An analytic argument and independent finite-chain calculation show why this witness cannot support the stated assertion. The supported Major impact is confined to this core negative derivation; the positive-curvature optimality theorem is not refuted, and substantial unverified appendix material is explicitly excluded. Scientific adjudication remains pending.\par
\section*{Paper identity and reviewed version}
This report concerns On Linear Convergence in Smooth Convex-Concave Bilinearly-Coupled Saddle-Point Optimization: Lower Bounds and Optimal Algorithms, by Ekaterina Borodich et al., ICML 2025, PMLR volume 267; identifier `I25-43671`, OpenReview `vVTgnjpaLp`. The reviewed object is the frozen published proceedings PDF, retrieved on 25 October 2026. Page references below are one-based PDF pages, which coincide with the internal page numbers. They are not the continuous PMLR volume pagination.\par
\section*{Scope and assumptions}
The review follows the finite-dimensional formulation, Assumptions 2.3–2.7, the oracle model in Definition 3.1, the negative result in Theorem 3.2 and Appendix F, the solution-set argument in Appendix E, and the upper-bound reduction in Appendices L–N. Theorem 3.3 and the recursive algorithm’s guarantees were read at the statement/dependency level; the long adversarial constructions in Appendices H–J and the full induction in K.1.1 were not independently verified. This report is therefore not a certification of the whole 56-page analysis. Algorithm 1 was statically inspected; no author code was executed. There are no experimental plots relevant to the finding. The exact theorem/proof pages 5 and 18 and the contraction calculation on page 55 were visually checked.\par
The problem is \(\min_x\max_y f(x)+\langle y,Bx\rangle-g(y)\) on finite-dimensional Euclidean spaces. Both functions are convex and smooth, with stated lower curvature bounds \(\mu_x,\mu_y\geq0\). Assumption 2.5 bounds the norm and appropriate positive singular values of \(B\), conditional on gradient-range inclusions. The effective curvatures are \[\delta_x=\mu_x+\mu_{xy}^2/L_y,\qquad \delta_y=\mu_y+\mu_{yx}^2/L_x.\] A declared zero lower curvature does not require the function’s largest valid strong-convexity modulus to be zero. This distinction prevents an easy strongly convex example from refuting a worst-case class statement.\par
\section*{The all-time lower-bound witness}
Problem 1 (Appendix F cannot establish the fixed-instance nonconvergence assertion).  Theorem 3.2, page 5, puts the existence of one problem before the universal quantifier over execution time, and claims \(\operatorname{dist}^2(x_o(\tau);S_x)>R_x^2/8\) for every admissible first-order method and every \(\tau>0\). The fixed finite-dimensional instance used in Appendix F, page 18, cannot have this property.\par
\subsection*{Exact dependency and source construction}
Equation reference: eq:family. The theorem assumes Assumptions 2.3–2.5 and \(\delta_x=0\), with arbitrary fixed \(R_x>0\) and accuracy \(\epsilon>0\). It is distinct from the positive-curvature lower bound in Theorem 3.3. Appendix F proposes \[\min_{u,v}\max_y h(u)+\frac{L_x}{2}\|v\|^2+\mu_{yx}\langle y,v\rangle-\frac{\mu_y}{2}\|y\|^2,\] where \(h\) is a finite-dimensional convex \(L_x\)-smooth hard function cited from Nesterov. It explicitly claims the solution is \((u^*,0,0)\) with \(u^*=\arg\min h\), then infers the entire theorem. The issue is the inference to an all-time distance obstruction, not the primal-dual gap comparison in this construction.\par
\subsection*{A general analytic check of this family}
Equation reference: eq:gd. Fix any dimension and any \(h\) as in the source’s proposed construction, with a minimizer. Starting at zero, keep \(v_k=y_k=0\) and update \[u_{k+1}=u_k-L_x^{-1}\nabla h(u_k).\] The full primal gradient at \((u_k,0)\) is \((\nabla h(u_k),0)\), so this is allowed by the gradient-span memory in Definition 3.1 and takes a finite time per iteration. No hidden access to \(u^*\) or non-oracle operation is used.\par
Equation reference: eq:fejer. For a convex \(L_x\)-smooth function, cocoercivity with a minimizer gives \(\langle\nabla h(u),u-u^*\rangle\geq L_x^{-1}\|\nabla h(u)\|^2\). This inequality follows by applying the smoothness descent inequality to \(h\) after subtracting the affine tangent at each of two points, then adding the resulting inequalities. Expansion of [eq:gd] (\#eq:gd) yields \[\|u_{k+1}-u^*\|^2\leq\|u_k-u^*\|^2-L_x^{-2}\|\nabla h(u_k)\|^2.\] Consequently the iterates are bounded and their gradients tend to zero. Finite dimensionality provides a convergent subsequence; continuity makes its limit a minimizer. Applying the same Fejer inequality to that minimizer forces the whole sequence to converge to it. Thus \(\operatorname{dist}(u_k,\arg\min h)\to0\). In particular, for the unique minimizer claimed in Appendix F, \(\|(u_k,0)-(u^*,0)\|^2<R_x^2/8\) at some finite time. This proves that none of the fixed solvable instances of the displayed family can establish the advertised universal-in-time lower bound.\par
\subsection*{Finite diagnostic and verification of assumptions}
An independent check makes the quantifier distinction concrete. Let \(u\in\mathbb R^{20}\), \(v,y\in\mathbb R\), let \(H\) have diagonal \(2\) and adjacent off-diagonals \(-1\), and set \(h(u)=u^THu/2-c e_1^Tu\), with \(c=\|H^{-1}e_1\|^{-1}\). Take \[(L_x,L_y,L_{xy},\mu_x,\mu_y,\mu_{xy},\mu_{yx})=(4,2,2,0,1,0,1).\] Then \(B=(0,\ldots,0,1)\), \(f(u,v)=h(u)+2v^2\), and \(g(y)=y^2/2\). The eigenvalues of \(H\) are \(2-2\cos(j\pi/21)\) and lie in \((0,4)\); the full primal Hessian is bounded by \(4I\). The dual Hessian is \(1\), so its declared strong convexity and smoothness bounds \(1\) and \(2\) are valid. The minimum eigenvalue of \(B^TB\) is zero, its positive eigenvalue and that of \(BB^T\) are \(1\), and \(\nabla g(y)\) lies in \(\operatorname{range}B=\mathbb R\). Thus Assumptions 2.3–2.5 hold, the extra equality condition when both coupling lower bounds are positive is inapplicable, and \(\delta_x=0\). The unique saddle is \((H^{-1}ce_1,0,0)\) and \(R_x=1\) by construction.\par
The independently written 20-dimensional gradient calculation crosses the claimed \(1/8\) barrier at iteration 148: squared distance \(0.1243643237\). At 1,000 iterations it is \(8.9039\times10^{-6}\). These figures verify the illustration, not a disproof by a single easy instance of an existential worst-case theorem. The decisive evidence is the preceding convergence proof for the whole proposed fixed-instance family.\par
\subsection*{Alternative interpretation and bounded consequence}
The intended lower-bound reading is plausibly horizon dependent: for each prescribed number of oracle calls, choose a sufficiently large dimension and a corresponding hard function. That allows class-wide slow convergence while every fixed finite-dimensional function eventually becomes solvable. It does not imply the source’s displayed order “one instance, all times.” Reinterpreting the proof this way changes a substantive quantifier. The later accuracy-dependent gap bound and the lack of a uniform linear rate may survive, but a complete replacement theorem must specify the dimension/horizon dependence and rederive the oracle-time relation. This review has not proved that replacement in all parameter regimes.\par
The supported assessment is Major for the invalid supporting derivation of the all-time negative claim. Section 2.2 explicitly uses Theorem 3.2 to present a necessary-and-sufficient characterization, making this a technical contribution rather than a peripheral display error. The report does not claim to have disproved every possible existential formulation under only Assumptions 2.3–2.5, especially formulations allowing unattained dual optima. Nor does it infer failure of the optimal positive-curvature complexity theorem. Appendix F’s own solvable witness is sufficient to establish the reported proof defect. No correction or later version is present among the frozen inputs; no exhaustive erratum search is asserted.\par
\section*{Checks that survive and review limits}
For the regularized upper-bound reduction, the coefficient arithmetic in Appendix M is consistent. The visually checked Equation (183) uses coefficients twelve, not one half, in the Lyapunov quantity: \[R_{\rm out}^2+12D_{f,\rm out}+12D_{g,\rm out}
\leq16\epsilon n R_{\rm in}^2+8D_{f,\rm in}+8D_{g,\rm in}.\] With \(\epsilon=1/(24n)\), both coefficients are exactly \(2/3\) of the corresponding input Lyapunov coefficients. The weighted operator bound \(\|P^{-1/2}Q_3P^{-1/2}\|=\|B\|/\sqrt{\delta_x\delta_y}\) is also consistent with \(\sqrt{\kappa_{xy}}\). Appendix L’s ascending ordering makes the ceiling-product estimates valid because each successive ratio is at least one. These checks neither establish the omitted induction in Appendix K nor independently validate all hard instances in I–J.\par
\section*{Conclusion}
Dimension | State\par
Supported assessment | Major: Appendix F does not prove the fixed-instance, all-time distance lower bound.\par
Repair verification | Horizon-dependent reading identified; complete corrected theorem unverified.\par
Surviving scope | Positive-curvature optimality theorem not refuted; selected upper-bound algebra checked.\par
Author adoption | Unknown.\par
Adjudication | Pending; focused independent AI acceptance completed; no human review.\par
Remaining work | Correct and prove quantifiers; independently verify long lower-bound constructions and recursive induction.\par
1 Ekaterina Borodich et al., On Linear Convergence in Smooth Convex-Concave Bilinearly-Coupled Saddle-Point Optimization: Lower Bounds and Optimal Algorithms, PMLR 267 (2025), published proceedings version. https://proceedings.mlr.press/v267/borodich25a.html.\par
\section*{Scope and limitations}
Source code supplied by the authors was not executed. The calculation is an independent finite or analytic diagnostic, not a reproduction of the original training experiments. A focused check revisited the stated Problems against the frozen original source, assumptions, decisive arguments, counterevidence and impact boundaries. No human review has occurred; scientific adjudication is pending and author adoption is unknown.\par
\end{document}
